A direct sequential scan of the concatenation through $3455$ finds its first occurrence at the boundary $3453\mid3454\mid3455$: the substring starts at the third digit of $3454$, continues through its fourth digit, then uses the first three digits of $3455$. It reads $54\,345$. The five immediately preceding digits are the last three digits of $3453$ followed by the first two digits of $3454$, namely $\boxed{45334}$. To verify first occurrence exhaustively, start with an empty string and append the decimal representation of each integer from $1$ to $3455$ in turn, checking every new five-digit window. No earlier window equals $54345$; the first begins at zero-based character index $12707$. This finite scan includes windows within numbers and across all digit-length boundaries.
